\chapter{Maximum Flow}

Picture a network of pipes. Water enters at one special place, sloshes through
a tangle of junctions and pipes of various widths, and leaves at another
special place. Each pipe has a limit -- a maximum rate it can carry before it
bursts. The question of this chapter is embarrassingly simple to ask and
genuinely deep to answer: \emph{how much water per second can you push from the
entrance to the exit?}

That is the maximum-flow problem. It does not look like sorting, or shortest
paths, or anything you have met so far, and the first honest reaction to it is
``surely I just send as much as I can down every pipe.'' But the pipes share
junctions, and water that piles up at a junction has nowhere to go -- whatever
flows in must flow out. The constraints interlock. Greedily flooding one path
can block a better global arrangement, and the whole charm of the subject is
learning how to \emph{undo} a greedy mistake without starting over.

By the end of this chapter you will be able to compute the maximum flow of any
network by hand and in code; you will understand the single most beautiful
theorem in the area -- that the most you can \emph{push} equals the cheapest
way to \emph{cut} the network in two -- and you will see how a problem about
water pipes secretly solves a problem about matchmaking, with no water in sight.

\begin{intuitionbox}
Here is the mental picture to carry the whole way through. Flow is not just
about where water \emph{can} go; it is about leaving yourself the freedom to
\emph{reroute} water you already committed. The central trick -- the residual
network -- is a bookkeeping device that says, ``if I previously sent five units
down this pipe, I am allowed to cancel up to five of them later by pretending
to send flow backward.'' That permission to take back a decision is what turns
a hopeless greedy idea into an exact algorithm.
\end{intuitionbox}

\section{Flow Networks}
\label{ch26:networks}

Let us pin down the objects. A \emph{flow network} is a directed graph
$G = (V, E)$ in which every edge $(u,v) \in E$ carries a nonnegative
\emph{capacity} $c(u,v) \ge 0$. Two vertices are singled out: a \emph{source}
$s$, where material is produced, and a \emph{sink} $t$, where it is consumed.
We assume every vertex lies on some path from $s$ to $t$, so the graph is
connected in the relevant sense and $\abs{E} \ge \abs{V} - 1$.

For convenience we adopt a small convention: if $(u,v)$ is not an edge, we set
$c(u,v) = 0$, and we forbid \emph{antiparallel} edges, meaning we never have
both $(u,v)$ and $(v,u)$ in $E$ at once. (If a real problem has both, split one
of them by inserting a dummy vertex in the middle; nothing important changes.
We will explain why this convention is convenient when residual networks
appear.)

\begin{definition}[Flow]
\label{ch26:def-flow}
A \emph{flow} in $G$ is a function $f : V \times V \to \R$ satisfying two
properties:
\begin{itemize}
  \item \textbf{Capacity constraint:} for all $u,v \in V$,
        $\;0 \le f(u,v) \le c(u,v)$.
  \item \textbf{Flow conservation:} for every vertex $u \in V \setminus
        \{s,t\}$,
        \[
          \sum_{v \in V} f(v,u) \;=\; \sum_{v \in V} f(u,v).
        \]
\end{itemize}
The \emph{value} of the flow, written $\abs{f}$, is the net amount leaving the
source,
\[
  \abs{f} \;=\; \sum_{v \in V} f(s,v) \;-\; \sum_{v \in V} f(v,s).
\]
\end{definition}

Read the two constraints slowly, because the entire chapter is a meditation on
them. The capacity constraint says no pipe carries more than its width allows,
and (in this nonnegative formulation) no pipe carries a negative amount. Flow
conservation is the law of the junction: at every vertex other than the source
and sink, total inflow equals total outflow. Water is neither created nor
destroyed inside the network; it only gathers at $s$ and drains at $t$.

\begin{focusbox}
A flow is a \emph{global} object that must be \emph{locally consistent}
everywhere. You cannot just maximize one pipe at a time, because changing one
pipe's flow ripples through conservation at both of its endpoints. The skill
you are building is reasoning about a quantity ($\abs{f}$) that lives at the
source and sink, while honoring constraints that live on every edge and at
every interior vertex.
\end{focusbox}

\subsection{A network to keep on the table}
\label{ch26:running-example}

We will trace one small network through the whole chapter. It has six vertices
$s, a, b, c, d, t$ and the capacities drawn below. Memorize nothing; just keep
this picture handy.

\begin{center}
\begin{tikzpicture}[
    >={Stealth[length=2.4mm]},
    node distance=18mm and 26mm,
    vtx/.style={circle,draw,minimum size=8mm,inner sep=0pt,font=\small},
    every edge/.style={draw,->,font=\scriptsize}
  ]
  \node[vtx,fill=sourceblue!18] (s) {$s$};
  \node[vtx,above right=of s] (a) {$a$};
  \node[vtx,below right=of s] (b) {$b$};
  \node[vtx,right=30mm of a] (c) {$c$};
  \node[vtx,right=30mm of b] (d) {$d$};
  \node[vtx,fill=targetgreen!22,below right=of c] (t) {$t$};
  \path
    (s) edge node[above left]{$10$} (a)
    (s) edge node[below left]{$10$} (b)
    (a) edge node[left]{$2$} (b)
    (a) edge node[above]{$8$} (c)
    (b) edge node[below]{$9$} (d)
    (c) edge node[right]{$6$} (d)
    (c) edge node[above right]{$10$} (t)
    (d) edge node[below right]{$10$} (t);
\end{tikzpicture}
\end{center}

Take a guess right now, before any theory: what do you think the maximum flow
from $s$ to $t$ is? Many people eyeball $s$'s two outgoing pipes, see capacity
$10 + 10 = 20$, and guess 20. Hold that guess. We will see it is too
optimistic, and the reason why is the heart of the subject.

\begin{aestheticbox}
Flow networks are not a toy. The same model computes how much data a computer
network can route between two machines, how many trucks a road system can carry
between two cities, how to schedule airline crews, how to segment an image into
foreground and background, how to assign students to schools, and how baseball
teams get mathematically eliminated from a pennant race. Whenever a resource
must travel through a shared, capacity-limited structure, maximum flow is
lurking. It is one of the highest-leverage algorithms an engineer can know,
precisely because so many problems quietly reduce to it.
\end{aestheticbox}

\section{The Ford--Fulkerson Method}
\label{ch26:ff}

We attack maximum flow with an idea so natural it is almost dumb: \emph{start
with zero flow, then keep finding a path from $s$ to $t$ along which you can
push a little more, and push it.} A path with spare room all the way from
source to sink is called an \emph{augmenting path}. As long as one exists, you
are not done; when none exists, the claim (to be proved) is that you have the
maximum.

This is called a \emph{method} rather than an algorithm because it leaves one
choice unspecified: \emph{which} augmenting path do you pick when several
exist? Different answers give different algorithms with wildly different running
times. We will nail down a good choice in Section~\ref{ch26:edmonds-karp}. For
now, focus on the engine.

\subsection{Why naive greedy is not enough}

Let us try the dumb version on a tiny three-edge graph and watch it fail, so
that the cure feels earned. Consider this network, with all capacities equal
to 1 except where noted.

\begin{center}
\begin{tikzpicture}[
    >={Stealth[length=2.4mm]},
    vtx/.style={circle,draw,minimum size=8mm,inner sep=0pt,font=\small},
    every edge/.style={draw,->,font=\scriptsize}
  ]
  \node[vtx,fill=sourceblue!18] (s) {$s$};
  \node[vtx,above right=10mm and 22mm of s] (u) {$u$};
  \node[vtx,below right=10mm and 22mm of s] (v) {$v$};
  \node[vtx,fill=targetgreen!22,below right=10mm and 22mm of u] (t) {$t$};
  \path
    (s) edge node[above left]{$1$} (u)
    (s) edge node[below left]{$1$} (v)
    (u) edge node[right]{$1$} (v)
    (u) edge node[above right]{$1$} (t)
    (v) edge node[below right]{$1$} (t);
\end{tikzpicture}
\end{center}

Suppose your first greedy path is $s \to u \to v \to t$, and you push 1 unit
along it. Now $s \to u$, $u \to v$, and $v \to t$ are all saturated. Look for
another augmenting path with spare capacity: $s \to v$ has room, but $v \to t$
is full; $u \to t$ has room, but $s \to u$ is full. There is no path of
strictly positive spare capacity from $s$ to $t$. Naive greedy stops, declaring
the answer 1.

But the true answer is 2: push 1 unit along $s \to u \to t$ and 1 unit along
$s \to v \to t$, using no middle edge at all. The greedy mistake was committing
flow to the middle edge $u \to v$. We need a mechanism to \emph{retract} that
unit -- to send flow ``backward'' along $u \to v$ to cancel it, freeing $u$ to
serve $t$ and $v$ to be served by $s$. That mechanism is the residual network,
and the figure below shows exactly the move it unlocks: a path that uses the
\emph{dashed backward copy} of the saturated middle edge.

\begin{figure}[H]
\centering
\begin{tikzpicture}[
    >={Stealth[length=2.2mm]},
    vtx/.style={circle,draw,minimum size=8mm,inner sep=0pt,font=\small},
    fwd/.style={draw,->,font=\scriptsize},
    bwd/.style={draw,->,dashed,trapred,font=\scriptsize}
  ]
  \node[vtx,fill=sourceblue!18] (s) {$s$};
  \node[vtx,above right=10mm and 24mm of s] (u) {$u$};
  \node[vtx,below right=10mm and 24mm of s] (v) {$v$};
  \node[vtx,fill=targetgreen!22,below right=10mm and 24mm of u] (t) {$t$};
  % forward residual edges (spare capacity): s-v and u-t are unused
  \path[fwd] (s) edge node[below left]{$1$} (v);
  \path[fwd] (u) edge node[above right]{$1$} (t);
  % backward residual edges (from the saturated s-u, u-v, v-t)
  \path[bwd] (u) edge node[above left]{$1$} (s);
  \path[bwd] (v) edge node[right]{$1$} (u);
  \path[bwd] (t) edge node[below right]{$1$} (v);
\end{tikzpicture}
\caption{The residual network $G_f$ after the greedy push of $1$ unit along
$s\to u\to v\to t$. Solid arrows are forward residual edges (spare capacity in a
real pipe); dashed red arrows are backward residual edges that grant permission
to cancel existing flow. The new augmenting path is $s\to v\to u\to t$: take
$s\to v$ (spare), then the backward edge $v\to u$ (which cancels the committed
middle unit), then $u\to t$ (spare). Pushing $1$ along it leaves the middle edge
empty and yields the optimal flow of value $2$.}
\label{ch26:fig-diamond-residual}
\end{figure}

\subsection{Residual networks}
\label{ch26:residual}

Given a flow $f$, the \emph{residual network} $G_f$ tells you, for every
ordered pair of vertices, how much \emph{additional} flow you may push there. It
has two kinds of edges, and the second kind is the clever one.

For each edge $(u,v) \in E$:
\begin{itemize}
  \item a \emph{forward residual edge} $(u,v)$ with residual capacity
        $c_f(u,v) = c(u,v) - f(u,v)$, the unused room in the pipe;
  \item a \emph{backward residual edge} $(v,u)$ with residual capacity
        $c_f(v,u) = f(u,v)$, equal to the flow currently in the pipe.
\end{itemize}
Edges of residual capacity $0$ are simply omitted. The backward edge encodes
exactly the ``permission to cancel'' from our intuition box: if $f(u,v)$ units
are flowing forward, you may push up to $f(u,v)$ units backward, which has the
net effect of reducing the forward flow.

\begin{formalbox}
Formally, the residual capacity is
\[
  c_f(u,v) =
  \begin{cases}
    c(u,v) - f(u,v) & \text{if } (u,v) \in E,\\
    f(v,u) & \text{if } (v,u) \in E,\\
    0 & \text{otherwise.}
  \end{cases}
\]
The residual network is $G_f = (V, E_f)$ where
$E_f = \{(u,v) \in V \times V : c_f(u,v) > 0\}$. Note $\abs{E_f} \le 2\abs{E}$,
so $G_f$ is no more than twice the size of $G$. This is exactly why we banned
antiparallel edges: it guarantees that a forward residual edge and a backward
residual edge never collide on the same ordered pair, so $c_f$ is
well-defined by a single formula.
\end{formalbox}

\begin{trapbox}
The backward edge is not a real pipe and carries no real water. It is a
\emph{bookkeeping} edge that exists only inside $G_f$. When an augmenting path
uses a backward residual edge $(v,u)$, you are not sending new water from $v$ to
$u$; you are \emph{reducing} the existing real flow $f(u,v)$. Mixing up
``residual flow'' with ``actual flow'' is the single most common confusion in
this chapter. Actual flow $f$ lives in $G$; residual capacities live in $G_f$;
the algorithm reads $G_f$ to decide how to update $f$.
\end{trapbox}

\subsection{Augmenting along a path}

An \emph{augmenting path} $p$ is any simple path from $s$ to $t$ in the residual
network $G_f$. Because every edge on it has $c_f > 0$, we can push flow along
the whole path. How much? Exactly the smallest residual capacity on the path --
the bottleneck:
\[
  c_f(p) \;=\; \min\{\, c_f(u,v) : (u,v) \text{ is on } p \,\}.
\]
To augment, we walk the path and, for each residual edge $(u,v)$, push $c_f(p)$
units. If $(u,v)$ is a forward edge of $G$, we \emph{increase} $f(u,v)$; if it
is a backward edge (meaning $(v,u) \in E$), we \emph{decrease} $f(v,u)$. In
symbols, the updated flow $f'$ satisfies, for each $(u,v)$ on $p$,
\[
  f'(u,v) = f(u,v) + c_f(p) \quad\text{(forward)},
  \qquad
  f'(v,u) = f(v,u) - c_f(p) \quad\text{(backward)}.
\]

\begin{lemma}[Augmentation produces a valid flow]
\label{ch26:lemma-valid}
If $f$ is a flow in $G$ and $p$ is an augmenting path in $G_f$, then $f'$
defined above is a flow in $G$, and $\abs{f'} = \abs{f} + c_f(p) > \abs{f}$.
\end{lemma}

\begin{proof}[Why this is true, gently]
Capacities first. On a forward edge we add at most $c_f(u,v) = c(u,v) -
f(u,v)$, so $f'(u,v) \le c(u,v)$, and since $c_f(p) > 0$ we keep $f'(u,v) \ge
0$. On a backward edge we subtract at most $c_f(v,u) = f(u,v)$, so $f'(u,v) \ge
0$, and it only decreased so it stays $\le c(u,v)$. Every edge remains legal.

Conservation next. Take any interior vertex $w$ on the path. The path enters
$w$ once and leaves $w$ once, so exactly two residual edges at $w$ change. Walk
through the four cases (the entering edge forward or backward, likewise the
leaving edge) and in each case the net inflow at $w$ changes by $+c_f(p)$ from
one edge and $-c_f(p)$ from the other -- conservation is preserved. (Pushing
$c_f(p)$ ``backward'' on a real edge reduces both that edge's contribution to
its head and to its tail by the same amount, which is exactly what keeps the
books balanced.) Vertices not on $p$ are untouched.

Finally the value. The path leaves $s$ along exactly one residual edge. If it
is a forward edge out of $s$, then $f(s,\cdot)$ increases by $c_f(p)$; if it is
a backward edge into $s$, then some $f(\cdot,s)$ decreases by $c_f(p)$. Either
way the net out-of-$s$ quantity $\abs{f}$ rises by exactly $c_f(p)$.
\end{proof}

This lemma is the workhorse. It says: \emph{find any $s$-$t$ path in $G_f$,
push its bottleneck, repeat.} Each augmentation strictly increases the value. We
state the method as pseudocode.

\begin{algorithm}[H]
\caption{\textsc{Ford-Fulkerson}$(G, s, t)$}
\KwIn{flow network $G$ with capacities $c$, source $s$, sink $t$}
\KwOut{a maximum flow $f$}
\For{each edge $(u,v) \in E$}{$f(u,v) \gets 0$}
\While{there exists an augmenting path $p$ from $s$ to $t$ in $G_f$}{
  $c_f(p) \gets \min\{ c_f(u,v) : (u,v) \in p\}$\;
  \For{each edge $(u,v)$ on $p$}{
    \uIf{$(u,v) \in E$}{$f(u,v) \gets f(u,v) + c_f(p)$}
    \Else{$f(v,u) \gets f(v,u) - c_f(p)$}
  }
}
\Return $f$\;
\end{algorithm}

\subsection{A full numeric trace}
\label{ch26:trace}

Let us run the method on the running network of
Section~\ref{ch26:running-example}. We start with $f = 0$ everywhere, so
initially $G_f$ is just $G$ with every edge at its full capacity.

\textbf{Augmentation 1.} Find a path; say we find $s \to a \to c \to t$. The
residual capacities are $10, 8, 10$, so the bottleneck is $c_f(p) = 8$. Push 8.
Now $f(s,a)=8$, $f(a,c)=8$, $f(c,t)=8$, and $\abs{f}=8$. The edge $a \to c$ is
saturated.

\textbf{Augmentation 2.} Find $s \to b \to d \to t$. Residual capacities
$10, 9, 10$; bottleneck $9$. Push 9. Now $f(s,b)=9$, $f(b,d)=9$, $f(d,t)=9$,
and $\abs{f}=8+9=17$. The edge $b \to d$ is saturated.

\textbf{Augmentation 3.} Is there any path left? Let us search the residual
network from $s$ and see how far we get. From $s$ we can still reach $a$
(residual $10-8=2$) and $b$ (residual $10-9=1$). From $a$ the forward edge
$a\to c$ is saturated, but $a\to b$ has residual $2$ -- so $a$ only takes us to
$b$, which we already have. From $b$ the lone out-edge $b\to d$ is saturated. So
the search from $s$ reaches exactly $\{s,a,b\}$ and stops dead.

You might hope a \emph{backward} residual edge rescues us, the way it did in the
diamond. The only real edge with spare capacity past the waist is $c\to d$
(residual $6$), which would let us try
\[
  s \;\to\; a \;\to\; c \;\to\; d \;\to\; t.
\]
But this path is blocked at its second step: $c_f(a,c)=0$, because $a\to c$ is
saturated. We can never \emph{reach} $c$ from $s$ in $G_f$, so its outgoing room
is useless. The honest conclusion: after pushing $8$ and $9$, \emph{there is no
augmenting path}, and the maximum flow of this network is
\[
  \abs{f} = 17,
\]
not the 20 you might have guessed from the source's two outgoing pipes of
capacity 10 each. The bottleneck is downstream: edges $a\to c$ (cap 8) and
$b \to d$ (cap 9) together can pass at most $8+9 = 17$, and everything from $s$
must funnel through them.

\begin{walkbox}
Go look at the running network again and find the ``narrowest waist'' of the
whole graph -- a set of edges you could cut to completely separate $s$ from $t$
whose total capacity is as small as possible. Did you find $\{a\to c,\,b\to
d\}$ with total $8 + 9 = 17$? Notice that this number is exactly the maximum
flow we computed, and that this waist is the very cut drawn in
Figure~\ref{ch26:fig-mincut} below. Sit with that coincidence on your walk. It
is not a coincidence; it is the deepest theorem in the chapter, coming up next.
\end{walkbox}

\section{Cuts and the Max-Flow Min-Cut Theorem}
\label{ch26:mincut}

The trace ended on a tantalizing note: the maximum flow equaled the capacity of
the cheapest ``waist'' separating $s$ from $t$. Let us make ``waist'' precise
and then prove the equality. This is the intellectual summit of the chapter.

\begin{definition}[Cut]
\label{ch26:def-cut}
An \emph{$s$-$t$ cut} of $G$ is a partition of $V$ into two sets $S$ and
$T = V \setminus S$ with $s \in S$ and $t \in T$. The \emph{net flow across the
cut} is
\[
  f(S,T) = \sum_{u \in S}\sum_{v \in T} f(u,v) - \sum_{u \in S}\sum_{v \in T} f(v,u),
\]
and the \emph{capacity of the cut} is
\[
  c(S,T) = \sum_{u \in S}\sum_{v \in T} c(u,v).
\]
\end{definition}

Notice the asymmetry: the \emph{net flow} subtracts the flow coming back from
$T$ to $S$, but the \emph{capacity} only counts edges going from $S$ to $T$
(forward), with no subtraction. That asymmetry is exactly where the slack
between flow and capacity will live.

\begin{lemma}[Flow across any cut equals the flow value]
\label{ch26:lemma-cut-value}
For any flow $f$ and any $s$-$t$ cut $(S,T)$, we have $f(S,T) = \abs{f}$.
\end{lemma}

\begin{proof}[The bookkeeping]
Start from $\abs{f} = \sum_v f(s,v) - \sum_v f(v,s)$, the net flow out of $s$
alone. Now \emph{add zero} cleverly: for every other vertex $u \in S$,
conservation says its net outflow $\sum_v f(u,v) - \sum_v f(v,u)$ is $0$, so we
may add all of these in without changing the total. Summing over all $u \in S$,
\[
  \abs{f} = \sum_{u \in S}\Bigl(\sum_{v \in V} f(u,v) - \sum_{v \in V} f(v,u)\Bigr).
\]
Split the inner sum over $v$ into $v \in S$ and $v \in T$. The $v \in S$ terms
cancel in pairs (each $f(u,v)$ with $u,v\in S$ appears once with a $+$ and
once, as $f(v,u)$, with a $-$). What survives is precisely the $S$-to-$T$ part,
which is $f(S,T)$. Done.
\end{proof}

This little identity has an enormous consequence. The flow value is the same no
matter which cut you measure it across -- pick any line separating $s$ from $t$,
and the net flow crossing it equals $\abs{f}$. Combined with capacities, it
gives an instant upper bound.

\begin{corollary}[Weak duality]
\label{ch26:weak-duality}
For any flow $f$ and any cut $(S,T)$, $\abs{f} \le c(S,T)$.
\end{corollary}

\begin{proof}
$\abs{f} = f(S,T) = \sum f(u,v) - \sum f(v,u) \le \sum f(u,v) \le \sum c(u,v)
= c(S,T)$, where the first inequality drops the nonnegative back-flow term and
the second uses the capacity constraint edge by edge.
\end{proof}

So \emph{every} flow value is at most \emph{every} cut capacity. In our running
example, the cut $S = \{s,a,b\}$, $T=\{c,d,t\}$ has forward edges $a\to c$ (cap
8) and $b\to d$ (cap 9), capacity $17$. Weak duality already guarantees the
maximum flow cannot exceed 17. The remarkable fact is that this bound is always
\emph{tight}: some flow actually achieves the smallest cut.

\begin{figure}[H]
\centering
\begin{tikzpicture}[
    >={Stealth[length=2.4mm]},
    node distance=18mm and 26mm,
    vtx/.style={circle,draw,minimum size=8mm,inner sep=0pt,font=\small},
    every edge/.style={draw,->,font=\scriptsize}
  ]
  \node[vtx,fill=sourceblue!18] (s) {$s$};
  \node[vtx,above right=of s] (a) {$a$};
  \node[vtx,below right=of s] (b) {$b$};
  \node[vtx,right=30mm of a] (c) {$c$};
  \node[vtx,right=30mm of b] (d) {$d$};
  \node[vtx,fill=targetgreen!22,below right=of c] (t) {$t$};
  \path
    (s) edge node[above left]{$10$} (a)
    (s) edge node[below left]{$10$} (b)
    (a) edge node[left]{$2$} (b)
    (a) edge[trapred,very thick] node[above]{$8$} (c)
    (b) edge[trapred,very thick] node[below]{$9$} (d)
    (c) edge node[right]{$6$} (d)
    (c) edge node[above right]{$10$} (t)
    (d) edge node[below right]{$10$} (t);
  % the cut line separating S={s,a,b} from T={c,d,t}
  \begin{scope}[on background layer]
    \draw[structurepurple,dashed,very thick]
      ($(a)!0.5!(c)+(0,9mm)$) -- ($(b)!0.5!(d)-(0,9mm)$);
    \node[structurepurple,font=\scriptsize] at ($(a)!0.5!(c)+(0,13mm)$) {cut};
  \end{scope}
  \node[font=\scriptsize,sourceblue!70!black] at ($(s)+(-2mm,-9mm)$) {$S=\{s,a,b\}$};
  \node[font=\scriptsize,targetgreen!60!black] at ($(t)+(2mm,9mm)$) {$T=\{c,d,t\}$};
\end{tikzpicture}
\caption{The minimum cut of the running network. The dashed purple line
separates $S=\{s,a,b\}$ from $T=\{c,d,t\}$. Only the two thick red edges $a\to
c$ and $b\to d$ cross forward from $S$ to $T$, so the cut capacity is $8+9=17$.
By max-flow min-cut this equals the maximum flow value -- the network's true
bottleneck lies at this downstream waist, not at the source's two pipes of
capacity $10$.}
\label{ch26:fig-mincut}
\end{figure}

\begin{theorem}[Max-Flow Min-Cut]
\label{ch26:mfmc}
For a flow $f$ in a network $G$ with source $s$ and sink $t$, the following
three statements are equivalent:
\begin{enumerate}
  \item $f$ is a maximum flow in $G$.
  \item The residual network $G_f$ contains no augmenting path from $s$ to $t$.
  \item $\abs{f} = c(S,T)$ for some $s$-$t$ cut $(S,T)$.
\end{enumerate}
Consequently the value of a maximum flow equals the capacity of a minimum cut.
\end{theorem}

\begin{proof}[The three-way cycle]
We show $(1)\Rightarrow(2) \Rightarrow (3) \Rightarrow (1)$.

\emph{$(1)\Rightarrow(2)$.} Contrapositive. If $G_f$ \emph{has} an augmenting
path $p$, then by Lemma~\ref{ch26:lemma-valid} we can push $c_f(p) > 0$ more
flow, so $f$ was not maximum.

\emph{$(2)\Rightarrow(3)$.} This is the magic step. Suppose $G_f$ has no
$s$-$t$ path. Define
\[
  S = \{\, v \in V : \text{there is a path from } s \text{ to } v \text{ in } G_f \,\},
  \qquad T = V \setminus S.
\]
Then $s \in S$ (trivially) and $t \in T$ (because no path reaches it), so
$(S,T)$ is a genuine $s$-$t$ cut. Now look at any real edge $(u,v)$ with $u \in
S$, $v \in T$. If it had spare capacity, $c_f(u,v) = c(u,v) - f(u,v) > 0$, then
$v$ would be reachable from $s$ in $G_f$, contradicting $v \in T$. So every
such edge is \emph{saturated}: $f(u,v) = c(u,v)$. Symmetrically, any real edge
$(v,u)$ with $v \in T$, $u \in S$ must carry \emph{zero} flow, for otherwise
the backward residual edge $(u,v)$ would have $c_f(u,v) = f(v,u) > 0$ and again
reach into $T$. Plug these into the cut identity:
\[
  \abs{f} = f(S,T) = \underbrace{\sum_{u\in S, v\in T} f(u,v)}_{=\,\sum c(u,v)\,=\,c(S,T)} - \underbrace{\sum_{u\in S, v\in T} f(v,u)}_{=\,0} = c(S,T).
\]

\emph{$(3)\Rightarrow(1)$.} If $\abs{f} = c(S,T)$ for some cut, then by weak
duality every flow $f'$ satisfies $\abs{f'} \le c(S,T) = \abs{f}$, so $f$ is
maximum. The final sentence follows because the maximum flow value equals this
particular cut capacity, and no cut can be cheaper (weak duality again), so it
is the minimum cut.
\end{proof}

\begin{focusbox}
The proof gives you a free gift: once Ford--Fulkerson halts, the set
$S$ of vertices reachable from $s$ in the final residual network \emph{is} a
minimum cut. So the same run that computes the maximum flow also hands you the
bottleneck -- the exact edges you would need to upgrade (or sabotage) to change
the answer. ``Run a max-flow, then read off the reachable set'' is a standard
move you will use again and again.
\end{focusbox}

\begin{historybox}
The Ford--Fulkerson method and the max-flow min-cut theorem were published in
1956 by Lester Ford and Delbert Fulkerson at the RAND Corporation, where the
motivating application was, of all things, estimating the capacity of the
Soviet railway network -- and, equivalently by min-cut, the cheapest way to
interdict it. The theorem is one of a famous family of min-max dualities in
combinatorics (alongside Konig's theorem and the duality of linear
programming, the subject of a later chapter). In CLRS this material is
Chapter~24, ``Maximum Flow'': Section~24.1 covers flow networks, 24.2 the
Ford--Fulkerson method and the max-flow min-cut theorem, 24.3 maximum bipartite
matching. (Older editions number it Chapter~26; if your copy does, look there.)
\end{historybox}

\subsection{Why ``method'' and not ``algorithm'': termination}
\label{ch26:termination}

We have been careful to call Ford--Fulkerson a \emph{method}, and now we can say
why precisely. If all capacities are integers, then every residual capacity
stays an integer, every bottleneck $c_f(p)$ is a positive integer, and so each
augmentation increases $\abs{f}$ by at least $1$. The value can rise at most up
to $\abs{f^*} = c(S,T)$ for the min cut, so the method halts after at most
$\abs{f^*}$ augmentations. Good -- it terminates, and as a bonus we get:

\begin{corollary}[Integrality theorem]
\label{ch26:integrality}
If every capacity is an integer, then some maximum flow is integer-valued
(every $f(u,v)$ is an integer). Ford--Fulkerson, started from $f=0$, produces
such a flow.
\end{corollary}

This integrality is not a footnote; it is the bridge to combinatorics. It is
why a problem about ``half a matching'' never arises -- if capacities are 1, the
flow on every edge is forced to be $0$ or $1$, which is precisely a yes/no
decision. We exploit this directly in Section~\ref{ch26:matching}.

\begin{trapbox}
With \emph{irrational} capacities, the bad choice of augmenting paths can make
Ford--Fulkerson run forever -- and worse, converge to a flow value strictly
\emph{below} the true maximum. There is a classic five-vertex example using the
golden ratio $\phi = (\sqrt5-1)/2$ where naive path selection produces an
infinite sequence of ever-smaller augmentations summing to a value well short
of optimal. The lesson is not ``avoid irrationals''; it is ``the path choice
matters.'' Choosing \emph{shortest} augmenting paths -- the Edmonds--Karp rule
-- cures this completely, regardless of the capacities.
\end{trapbox}

\section{Edmonds--Karp: Choosing Good Paths}
\label{ch26:edmonds-karp}

The fix is almost insultingly simple. When several augmenting paths exist,
always take a \emph{shortest} one -- fewest edges -- which you find with a
breadth-first search from $s$ in $G_f$. That single rule turns the method into a
polynomial-time algorithm whose running time does not depend on the capacity
values at all.

\begin{algorithm}[H]
\caption{\textsc{Edmonds-Karp}$(G,s,t)$}
\KwIn{flow network $G$, source $s$, sink $t$}
\KwOut{a maximum flow $f$}
$f \gets 0$ on every edge\;
\While{BFS from $s$ in $G_f$ reaches $t$}{
  let $p$ be the BFS tree path $s \rightsquigarrow t$ (fewest edges)\;
  $c_f(p) \gets \min\{c_f(u,v) : (u,v) \in p\}$\;
  augment $f$ along $p$ by $c_f(p)$ \tcp*{as in Ford--Fulkerson}
}
\Return $f$\;
\end{algorithm}

Why does ``shortest path'' help? The key is a monotonicity fact about BFS
distances in the residual network. Write $\delta_f(s,v)$ for the
\emph{breadth-first distance} (number of edges) from $s$ to $v$ in $G_f$.

\begin{lemma}[Distances never decrease]
\label{ch26:lemma-distances}
Running Edmonds--Karp, for every vertex $v \neq s,t$ the distance
$\delta_f(s,v)$ in the residual network is monotonically nondecreasing across
augmentations.
\end{lemma}

The proof (we sketch it rather than belabor it) is a tidy argument by
contradiction: suppose some augmentation strictly decreased $\delta(s,v)$, and
take the $v$ closest to $s$ for which this happens. Augmenting along a
\emph{shortest} path can only \emph{add} residual edges that go ``backward''
one BFS level, never a shortcut that jumps forward, so no distance can drop --
contradiction. The discipline of always using a shortest path is exactly what
forbids the creation of shortcuts.

\begin{lemma}[At most $O(VE)$ augmentations]
\label{ch26:lemma-count}
Edmonds--Karp performs $O(\abs{V}\abs{E})$ augmentations before halting.
\end{lemma}

\begin{proof}[The critical-edge counting argument]
Call an edge $(u,v)$ \emph{critical} on an augmentation if it is the bottleneck,
i.e.\ $c_f(u,v) = c_f(p)$; this edge disappears from $G_f$ (its residual
capacity drops to $0$). On a shortest augmenting path through $(u,v)$ we have
$\delta_f(s,v) = \delta_f(s,u) + 1$. For $(u,v)$ to become critical \emph{again}
later, the flow on it must first increase from $0$, which happens only when
$(u,v)$ lies on an augmenting path in the other direction, i.e.\ when $(v,u)$
is used; at that moment $\delta_{f'}(s,u) = \delta_{f'}(s,v) + 1$. Because
distances never decrease (Lemma~\ref{ch26:lemma-distances}), between the two
times $(u,v)$ is critical, $\delta(s,u)$ must have grown by at least $2$. Since
$\delta$ ranges in $\{0,1,\dots,\abs{V}\}$ (or is $\infty$), each edge can be
critical at most $\abs{V}/2 = O(\abs{V})$ times. There are $O(\abs{E})$
residual edges, so the total number of critical events -- and hence
augmentations, since each augmentation has at least one critical edge -- is
$O(\abs{V}\abs{E})$.
\end{proof}

Each augmentation costs one BFS, which is $O(\abs{V} + \abs{E}) = O(\abs{E})$
on a connected network. Multiply.

\begin{theorem}[Edmonds--Karp running time]
\label{ch26:ek-time}
Edmonds--Karp computes a maximum flow in $O(\abs{V}\abs{E}^2)$ time, regardless
of the edge capacities.
\end{theorem}

\begin{complexitybox}
\textbf{Ford--Fulkerson} (any augmenting path, integer capacities): up to
$\abs{f^*}$ augmentations, each $O(E)$, giving $O(E\cdot\abs{f^*})$. This is
\emph{pseudo-polynomial} -- it depends on the numeric value of the answer, so a
single capacity of $1{,}000{,}000$ can make it crawl.

\textbf{Edmonds--Karp} (shortest augmenting path via BFS):
$O(V E^2)$, independent of capacities.

For reference, more advanced algorithms beat this: Dinic's algorithm runs in
$O(V^2 E)$; push--relabel methods reach $O(V^2\sqrt{E})$ or $O(V^3)$; and modern
results (Orlin; and the 2022 almost-linear-time breakthrough) push the
frontier far lower. Edmonds--Karp is the right first algorithm to \emph{know
cold}, because it is Ford--Fulkerson plus one line: ``use BFS.''
\end{complexitybox}

\begin{labbox}
Re-run the running network of Section~\ref{ch26:running-example} \emph{by hand}
under the Edmonds--Karp rule, and keep a table.

\begin{center}
\begin{tabular}{c l c c}
\toprule
step & shortest path $p$ & bottleneck $c_f(p)$ & $\abs{f}$ after \\
\midrule
0 & --- & --- & $0$ \\
1 & $s\to a\to c\to t$ & $\min(10,8,10)=8$ & $8$ \\
2 & $s\to b\to d\to t$ & $\min(10,9,10)=9$ & $17$ \\
3 & no $s$-$t$ path in $G_f$ & --- & $17$ (max) \\
\bottomrule
\end{tabular}
\end{center}

Now do the verification the theorem promises: after step 2, run a BFS from $s$
in $G_f$ and list which vertices you can reach. You should reach $\{s,a,b\}$ and
no further ($a\to c$ and $b\to d$ are saturated, blocking the way). That set is
$S$; the cut $(S,T)$ has capacity $c(a,c)+c(b,d) = 8+9 = 17 = \abs{f}$. You have
just \emph{certified} optimality with a min cut, by hand. Try inventing a third
path that increases the flow; convince yourself it is impossible.
\end{labbox}

\section{Application: Maximum Bipartite Matching}
\label{ch26:matching}

Now the payoff that makes the whole machine feel magical. We will solve a
problem that has nothing to do with water -- matchmaking -- by quietly building
a flow network and pushing ones and zeros through it.

A \emph{bipartite graph} $G = (L \cup R, E)$ has its vertices split into two
sides $L$ and $R$, with every edge joining a vertex of $L$ to a vertex of $R$.
Think of $L$ as job applicants and $R$ as positions, with an edge $(\ell, r)$
meaning ``applicant $\ell$ is qualified for position $r$.'' A \emph{matching} is
a set of edges $M \subseteq E$ no two of which share a vertex: each applicant
gets at most one job, each job at most one applicant. The \emph{maximum
matching} problem asks for a matching of greatest size. How many people can be
hired at once?

\begin{intuitionbox}
A matching is just an honest pairing-off with no double-booking. The trouble is
local greed: hire applicant 1 for the only job she shares with applicant 2, and
you may strand applicant 2 with nothing -- when a smarter pairing could employ
both. This is the very same ``my early commitment blocks a better global plan''
disease we cured with residual networks. So we will let flow's
backward-edge magic do the un-committing for us.
\end{intuitionbox}

\subsection{The reduction}

Build a flow network $G' = (V', E')$ from the bipartite graph as follows.
\begin{itemize}
  \item Add a new source $s$ and sink $t$.
  \item Add an edge $s \to \ell$ of capacity $1$ for every $\ell \in L$.
  \item Add an edge $r \to t$ of capacity $1$ for every $r \in R$.
  \item For every original edge $(\ell, r) \in E$, add a directed edge
        $\ell \to r$ of capacity $1$ (any capacity $\ge 1$ works; we use 1).
\end{itemize}
Every capacity is $1$. Here is the picture, with three applicants and three
jobs.

\begin{center}
\begin{tikzpicture}[
    >={Stealth[length=2.2mm]},
    vtx/.style={circle,draw,minimum size=7mm,inner sep=0pt,font=\small},
    every edge/.style={draw,->,font=\scriptsize}
  ]
  \node[vtx,fill=sourceblue!18] (s) {$s$};
  \node[vtx,above right=6mm and 18mm of s] (l1) {$\ell_1$};
  \node[vtx,right=18mm of s] (l2) {$\ell_2$};
  \node[vtx,below right=6mm and 18mm of s] (l3) {$\ell_3$};
  \node[vtx,right=22mm of l1] (r1) {$r_1$};
  \node[vtx,right=22mm of l2] (r2) {$r_2$};
  \node[vtx,right=22mm of l3] (r3) {$r_3$};
  \node[vtx,fill=targetgreen!22,right=18mm of r2] (t) {$t$};
  \path
    (s) edge (l1) (s) edge (l2) (s) edge (l3)
    (l1) edge (r1) (l1) edge (r2)
    (l2) edge (r1)
    (l3) edge (r2) (l3) edge (r3)
    (r1) edge (t) (r2) edge (t) (r3) edge (t);
\end{tikzpicture}
\end{center}

\begin{theorem}[Matching equals integer flow]
\label{ch26:matching-thm}
The maximum number of edges in a matching of $G$ equals the value of a maximum
flow in $G'$. Moreover, given an integer maximum flow $f$, the set of middle
edges $\ell \to r$ carrying flow $1$ is a maximum matching.
\end{theorem}

\begin{proof}[Why the correspondence works, both directions]
\emph{Matching $\to$ flow.} Given a matching $M$, route one unit along
$s \to \ell \to r \to t$ for each edge $(\ell,r)\in M$. Because $M$ has no two
edges sharing an $\ell$, each $s\to\ell$ carries at most $1$; because no two
share an $r$, each $r\to t$ carries at most $1$; capacities are respected and
conservation holds. This is a valid flow of value $\abs{M}$.

\emph{Flow $\to$ matching.} By the integrality theorem
(Corollary~\ref{ch26:integrality}), a maximum flow may be taken integer-valued,
so every middle edge carries $0$ or $1$. Let $M$ be the middle edges carrying
$1$. Each $\ell$ has only one incoming unit (from $s$, capacity $1$), so by
conservation it sends out at most one unit -- $\ell$ touches at most one matched
edge. Each $r$ sends only one unit to $t$, so receives at most one -- $r$
touches at most one matched edge. Thus $M$ is a matching, and $\abs{M} =
\abs{f}$ because every unit out of $s$ passes through exactly one middle edge.

Since the two constructions preserve value and are inverse in size, the maxima
coincide.
\end{proof}

So to find a maximum matching, build $G'$ and run Edmonds--Karp. Because every
capacity is $1$, the analysis even sharpens: the flow value is at most
$\min(\abs{L},\abs{R}) \le \abs{V}$, so Ford--Fulkerson alone makes at most
$\abs{V}$ augmentations, each an $O(E)$ search, giving a clean $O(VE)$ bound for
bipartite matching.

\subsection{Tracing the matching}

Run it on the three-by-three example. First augmenting path $s\to\ell_1\to r_1
\to t$, push 1: $\ell_1$ matched to $r_1$. Second path $s\to\ell_3\to r_2\to t$,
push 1: $\ell_3$ matched to $r_2$. Now try $\ell_2$: from $s\to\ell_2$ we reach
$r_1$, but $r_1\to t$ is saturated. Dead end? No -- the residual network offers
the backward edge $r_1\to\ell_1$ (because $f(\ell_1,r_1)=1$). The BFS finds
\[
  s \to \ell_2 \to r_1 \xrightarrow{\text{backward}} \ell_1 \to r_2
      \xrightarrow{\text{backward}} \ell_3 \to r_3 \to t.
\]
Pushing 1 unit along it \emph{re-routes} the earlier matches: $\ell_1$ drops
$r_1$ and takes $r_2$; $\ell_3$ drops $r_2$ and takes $r_3$; and $\ell_2$ grabs
the freed $r_1$. Final flow value $3$ -- a \emph{perfect} matching
$\{\ell_1 r_2, \ell_2 r_1, \ell_3 r_3\}$ that naive greed would have missed.
Figure~\ref{ch26:fig-matching-aug} traces that zig-zag.

\begin{figure}[H]
\centering
\begin{tikzpicture}[
    >={Stealth[length=2.2mm]},
    vtx/.style={circle,draw,minimum size=7mm,inner sep=0pt,font=\small},
    matched/.style={draw,->,targetgreen,very thick},
    cancel/.style={draw,->,dashed,trapred,font=\scriptsize},
    add/.style={draw,->,sourceblue,thick}
  ]
  \node[vtx] (l1) {$\ell_1$};
  \node[vtx,below=8mm of l1] (l2) {$\ell_2$};
  \node[vtx,below=8mm of l2] (l3) {$\ell_3$};
  \node[vtx,right=32mm of l1] (r1) {$r_1$};
  \node[vtx,right=32mm of l2] (r2) {$r_2$};
  \node[vtx,right=32mm of l3] (r3) {$r_3$};
  % current matches before the augmentation (forward, matched), bent up
  \draw[matched] (l1) to[bend left=14] node[above,font=\scriptsize]{$\ell_1 r_1$} (r1);
  \draw[matched] (l3) to[bend left=14] node[below,font=\scriptsize]{$\ell_3 r_2$} (r2);
  % the augmenting (alternating) path
  \path[add]    (l2) edge node[pos=0.45,above,font=\scriptsize,sourceblue]{1} (r1);
  \draw[cancel] (r1) to[bend left=14] node[below,font=\scriptsize]{back} (l1);
  \path[add]    (l1) edge node[pos=0.4,above,font=\scriptsize,sourceblue]{3} (r2);
  \draw[cancel] (r2) to[bend left=14] node[above,font=\scriptsize]{back} (l3);
  \path[add]    (l3) edge node[pos=0.45,below,font=\scriptsize,sourceblue]{5} (r3);
\end{tikzpicture}
\caption{The augmenting (alternating) path that grows the matching from $2$ to
$3$. Thick green edges are the current matches $\ell_1 r_1$ and $\ell_3 r_2$.
The path enters along blue ``add'' edges ($\ell_2 r_1$, $\ell_1 r_2$, $\ell_3
r_3$, numbered in traversal order) and crosses two dashed red backward edges
that \emph{undo} the old matches. Flipping every edge on the path turns the two
matches into three: $\{\ell_1 r_2,\ \ell_2 r_1,\ \ell_3 r_3\}$.}
\label{ch26:fig-matching-aug}
\end{figure}

\begin{focusbox}
That backward-edge zig-zag has a classical name in matching theory: an
\emph{augmenting path} for matchings -- a path that alternates between unmatched
and matched edges, starting and ending at unmatched vertices. Flipping it
increases the matching by one. The flow algorithm rediscovers this idea for
free: matching's augmenting paths are exactly $s$-$t$ paths in the residual
network. Two subjects, one mechanism.
\end{focusbox}

\begin{aestheticbox}
This reduction is the reason maximum flow earns its place in every algorithms
course. With it you also get, almost for free: assigning taxis to riders,
scheduling exams to rooms, detecting whether a sports team is mathematically
eliminated, segmenting medical images (foreground/background as a min cut),
and the entire theory of bipartite assignment that underlies kidney-exchange
matching programs and ad auctions. The pattern -- ``model the decisions as
edges of capacity 1, route flow, read the integer answer'' -- is one of the
most reusable tools in applied computing.
\end{aestheticbox}

\section{What to Carry Forward}
\label{ch26:recap}

\begin{recoverybox}
Five sentences to keep:
\begin{enumerate}
  \item A flow obeys \emph{capacity} on every edge and \emph{conservation} at
        every interior vertex; its \emph{value} is the net amount leaving $s$.
  \item The \emph{residual network} adds backward edges that grant permission
        to cancel earlier flow; this is what lets a greedy method become exact.
  \item Ford--Fulkerson repeatedly pushes the bottleneck of an augmenting path;
        with integer capacities it terminates and yields an integer flow.
  \item \emph{Max-flow equals min-cut}: when no augmenting path remains, the
        set reachable from $s$ in $G_f$ is a minimum cut certifying optimality.
  \item \emph{Edmonds--Karp} = ``always use the shortest augmenting path
        (BFS),'' giving $O(VE^2)$ time, independent of capacities; bipartite
        matching reduces to a unit-capacity flow.
\end{enumerate}
\end{recoverybox}

\section*{Exercises}
\addcontentsline{toc}{section}{Exercises}

\begin{exercise}[Warmup: read a flow]
\label{ch26:ex-read}
In the running network of Section~\ref{ch26:running-example}, suppose someone
hands you the assignment $f(s,a)=8$, $f(a,c)=8$, $f(c,t)=8$, $f(s,b)=9$,
$f(b,d)=9$, $f(d,t)=9$, and $0$ on the remaining edges. Verify that it satisfies
the capacity constraint and flow conservation at $a,b,c,d$, and compute
$\abs{f}$.
\end{exercise}

\begin{exercise}[Build a residual network]
\label{ch26:ex-residual}
For the flow in Exercise~\ref{ch26:ex-read}, draw the residual network $G_f$:
list every residual edge and its residual capacity (remember both forward and
backward edges). Then argue from your drawing that no $s$-$t$ path exists.
\end{exercise}

\begin{exercise}[Find the min cut]
\label{ch26:ex-cut}
Continuing the previous exercise, compute the set $S$ of vertices reachable from
$s$ in $G_f$, write down the cut $(S,T)$, and verify that its capacity equals
$\abs{f}$. Which two edges form the bottleneck?
\end{exercise}

\begin{exercise}[Greedy can fail]
\label{ch26:ex-greedy}
Recall the five-vertex network of Section~\ref{ch26:ff} (the $s,u,v,t$ diamond
with the middle edge $u\to v$, all capacities $1$). Trace Ford--Fulkerson
starting with the bad first path $s\to u\to v\to t$. Show explicitly the
augmenting path through the backward residual edge that recovers the optimal
flow of value $2$, and give the resulting flow on each real edge.
\end{exercise}

\begin{exercise}[Why no antiparallel edges]
\label{ch26:ex-antiparallel}
Suppose a network genuinely has both $(u,v)$ and $(v,u)$ as edges with positive
capacities. Explain in one or two sentences why the single-formula definition of
$c_f$ in Section~\ref{ch26:residual} would become ambiguous. Then describe the
fix (split one edge with a dummy vertex $w$: replace $(u,v)$ by $u\to w\to v$,
each of capacity $c(u,v)$) and argue it does not change the maximum flow value.
\end{exercise}

\begin{exercise}[Integrality is essential for matching]
\label{ch26:ex-integrality}
Show by a small example that if you allowed \emph{fractional} flow on the middle
edges of the bipartite construction, a maximum flow could split a unit of an
applicant across two jobs. Explain why
Corollary~\ref{ch26:integrality} guarantees we can always avoid this, so the
flow really does encode an honest matching.
\end{exercise}

\begin{exercise}[Counting augmentations]
\label{ch26:ex-count}
Construct a network on which \emph{naive} Ford--Fulkerson, making the worst
possible path choices, performs about $2{,}000{,}000$ augmentations, while
Edmonds--Karp performs only a handful. \emph{(Hint: the classic ``diamond''
with two outer edges of capacity $1{,}000{,}000$ and a middle edge of capacity
$1$; alternately augment through the middle.)} What does this say about the
difference between pseudo-polynomial and polynomial time?
\end{exercise}

\begin{exercise}[Vertex capacities]
\label{ch26:ex-vertex-cap}
Some problems limit how much flow can pass \emph{through a vertex} $v$, not just
along edges -- say at most $\mathrm{cap}(v)$ units. Show how to model this with
ordinary edge capacities by splitting $v$ into $v_{\text{in}}$ and
$v_{\text{out}}$ joined by an edge of capacity $\mathrm{cap}(v)$, routing all
in-edges to $v_{\text{in}}$ and all out-edges from $v_{\text{out}}$. Argue
correctness.
\end{exercise}

\begin{exercise}[Multiple sources and sinks]
\label{ch26:ex-multi}
A factory network has three sources $s_1,s_2,s_3$ (each able to produce any
amount) and two sinks $t_1,t_2$. Explain how to compute the maximum total flow
from all sources to all sinks by adding a single \emph{super-source} and a
single \emph{super-sink} with appropriate infinite-capacity edges, reducing to
the single-source single-sink problem of this chapter.
\end{exercise}

\begin{exercise}[Edge-disjoint paths -- a challenge]
\label{ch26:ex-disjoint}
Menger's theorem states that the maximum number of \emph{edge-disjoint} paths
from $s$ to $t$ in a directed graph equals the minimum number of edges whose
removal disconnects $t$ from $s$. Show how to prove this using max-flow
min-cut: assign every edge capacity $1$, and interpret an integer flow of value
$k$ as $k$ edge-disjoint paths. Where exactly does the integrality theorem enter
the argument?
\end{exercise}

\begin{exercise}[Uniqueness -- a challenge]
\label{ch26:ex-unique}
The \emph{value} of the maximum flow is unique, but the flow itself and the
minimum cut need not be. Construct a small network that has (a) two different
maximum flows of the same value, and (b) two different minimum cuts of the same
capacity. Then prove the general fact that if $(S_1,T_1)$ and $(S_2,T_2)$ are
both minimum cuts, so are $(S_1\cap S_2,\,\cdot)$ and $(S_1\cup S_2,\,\cdot)$ --
the minimum cuts form a lattice.
\end{exercise}
